% !TEX root = ../metatime_monograph.tex
\chapter{Decuplet Coefficients and Masses}
\label{ch:decuplet_predictions}

\section{Decuplet GMO Coefficients}

The decuplet follows an equal-spacing rule in hypercharge \(Y\):

\begin{equation}
M = M_0' + \alpha_{10} + \beta_{10} Y.
\label{eq:decuplet_formula}
\end{equation}

The coefficients are derived from the L-constants:

\subsection{Coefficient \(\alpha_{10}\)}

\begin{equation}
\alpha_{10} = E_{\text{proton}}\,\kappa\,
\frac{L_3^2 + L_4^2 + L_5^2 + L_3 + L_4 + L_5 + L_4}{L_4}.
\label{eq:alpha10}
\end{equation}

\subsubsection{Term-by-Term Breakdown}

\[
\begin{aligned}
L_3^2 &= 49,\\
L_4^2 &= 4,\\
L_5^2 &= 25,\\
L_3   &= 7,\\
L_4   &= 2,\\
L_5   &= 5,\\
L_4   &= 2.
\end{aligned}
\]

Sum: \(49 + 4 + 25 + 7 + 2 + 5 + 2 = 94\).

Dividing by \(L_4 = 2\): \(94/2 = 47\).

\subsubsection{Numerical Evaluation}

\[
\begin{aligned}
\alpha_{10} &= 938.272 \times 0.00919315 \times 47 \\
            &= 8.62573 \times 47 \\
            &= 405.41\ \mev.
\end{aligned}
\]

\subsection{Coefficient \(\beta_{10}\)}

\begin{equation}
\beta_{10} = -E_{\text{proton}}\,\kappa\,
\frac{L_3 L_5}{L_4}.
\label{eq:beta10}
\end{equation}

\subsubsection{Numerical Evaluation}

\[
\begin{aligned}
\frac{L_3 L_5}{L_4} &= \frac{7 \times 5}{2} = \frac{35}{2} = 17.5, \\[4pt]
\beta_{10} &= -938.272 \times 0.00919315 \times 17.5 \\
           &= -8.62573 \times 17.5 \\
           &= -150.95\ \mev.
\end{aligned}
\]

\section{Mass Calculation}

\[
\begin{aligned}
M_{\Delta^{++}} &= M_0' + \alpha_{10} + \beta_{10}(+1) \\
                &= 972.72 + 405.41 - 150.95 \\
                &= 1227.18\ \mev, \\[4pt]
M_{\Sigma^{*+}} &= M_0' + \alpha_{10} + \beta_{10}(0) \\
                &= 972.72 + 405.41 \\
                &= 1378.13\ \mev, \\[4pt]
M_{\Xi^{*0}} &= M_0' + \alpha_{10} + \beta_{10}(-1) \\
             &= 972.72 + 405.41 + 150.95 \\
             &= 1529.08\ \mev, \\[4pt]
M_{\Omega^-} &= M_0' + \alpha_{10} + \beta_{10}(-2) \\
             &= 972.72 + 405.41 + 301.90 \\
             &= 1680.03\ \mev.
\end{aligned}
\]

\section{Complete Table}

\[
\begin{array}{lccc}
\text{Particle} & \text{Predicted (MeV)} & \text{PDG (MeV)} & \text{Error} \\
\hline
\Delta^{++} & 1227.18 & 1232.0  & -0.39\% \\
\Sigma^{*+} & 1378.13 & 1382.8  & -0.34\% \\
\Xi^{*0}    & 1529.08 & 1531.8  & -0.18\% \\
\Omega^-   & 1680.03 & 1672.45 & +0.45\%
\end{array}
\]

Mean absolute error: \(0.34\%\).

\section{Discussion}

The decuplet fit is excellent: \(0.34\%\) mean error with no free parameters.
The equal-spacing rule is satisfied to high precision.  The systematic
pattern (negative error for lighter states, positive for \(\Omega^-\))
suggests that the true equal spacing is slightly smaller than our
\(\beta_{10} = -150.95\)~MeV, but the deviation is within 0.5\%.

The decuplet is thus one of the best-described sectors in the Metatime
framework.
